% Written by travel_notes 0.2.0 from a Markdown note.
\documentclass{amsart}
\usepackage[T1]{fontenc}
\usepackage{amssymb}
\usepackage[normalem]{ulem}
\usepackage{hyperref}

\newtheorem{theorem}{Theorem}
\newtheorem{lemma}[theorem]{Lemma}
\newtheorem{proposition}[theorem]{Proposition}
\newtheorem{corollary}[theorem]{Corollary}
\newtheorem{conjecture}[theorem]{Conjecture}
\newtheorem{claim}[theorem]{Claim}
\theoremstyle{definition}
\newtheorem{definition}[theorem]{Definition}
\newtheorem{example}[theorem]{Example}
\newtheorem{exercise}[theorem]{Exercise}
\theoremstyle{remark}
\newtheorem{remark}[theorem]{Remark}

\begin{document}

\title{Notes on isogeny graphs}
\maketitle

These notes show headings with ids, links to them, theorem-like blocks and comments in
math. See \hyperref[tn-h-a-proof]{the proof} and \hyperref[tn-h-remarks-questions]{the last section}.

\section{Definitions}\label{tn-h-definitions}

\begin{definition}
Let \(\ell \ne p\) be primes. The graph \(G_p(\ell)\) has the supersingular
\(j\)-invariants over \(\overline{\mathbb{F}}_p\) as vertices and \(\ell\)-isogenies as edges.

\end{definition}

\begin{theorem}[{Pizer}]
The graph \(G_p(\ell)\) is Ramanujan:
\[
|\lambda| \le 2\sqrt{\ell}
\]
for every eigenvalue \(\lambda \ne \ell + 1\).

\end{theorem}

\begin{lemma}
A walk of length \(O(\log p)\) mixes.

\end{lemma}

\section{A proof}\label{tn-h-a-proof}

\begin{proof}[{Proof (of the lemma)}]
Expand in an eigenbasis.

\begin{itemize}
\item The trivial eigenvalue gives the uniform distribution.
\item The others decay like \((2\sqrt{\ell}/(\ell+1))^k\).
\end{itemize}

\end{proof}

\begin{remark}
an ordinary quotation follows.

\end{remark}

\begin{quote}
Not a theorem: \textbf{Theorem.} appears later in the line.

\end{quote}

\begin{quote}
\textbf{Theorem 7.} A number written by hand is left alone.

\end{quote}

\begin{theorem}
in lower case, without a full stop, still counts.

\end{theorem}

\section{Comments in math}\label{tn-h-comments-in-math}

\newcommand{\Hom}{\operatorname{Hom}}
\begin{align}
\Hom(E, E') &\cong I \label{eq:hom} \\ \deg \varphi &= 100\% \cdot n(I)
- 0 \\
\# S &= 1
+ 2
\end{align}

Inline too: \(a + b c\), and a percent sign kept: \(50\%\).

\section{A proof}\label{tn-h-a-proof-2}

The second heading with this text gets its own id. So does a heading with no letters:

\section{2024}\label{tn-h-2024}

\subsection{\emph{Remarks} \& \texttt{questions}}\label{tn-h-remarks-questions}

\begin{enumerate}
\item[$\boxtimes$] a finished task
\item[$\square$] an open one, \hyperref[tn-h-2024]{linked}
\end{enumerate}

Struck \sout{through} text. Text with LaTeX's special characters: 100\% of \$5 \& \#1 \{x\} a\_b x\textasciicircum{}2 \textasciitilde{} \textbackslash{} \textless{} \textgreater{} \textbar{} \textquotedbl{}q\textquotedbl{}.
A line that starts with a bracket after a break:\\
{[}bracket] and \texttt{code with \textbackslash{} and \{braces\}}.

{\ttfamily\noindent
\mbox{}\textbackslash{}begin\{verbatim\}\\
\mbox{}nested\\
\mbox{}\textbackslash{}end\{verbatim\}\\
\par}

\end{document}
